\section{Relative effective Cartier divisors}
\label{section-effective-Cartier-morphisms}

\noindent
The following lemma shows that an effective Cartier divisor which is
flat over the base is really a ``family of effective Cartier divisors''
over the base. For example the restriction to any fibre is an effective
Cartier divisor.

\begin{lemma}
\label{lemma-relative-Cartier}
Let $f : X \to S$ be a morphism of schemes.
Let $D \subset X$ be a closed subscheme.
Assume
\begin{enumerate}
\item $D$ is an effective Cartier divisor, and
\item $D \to S$ is a flat morphism.
\end{enumerate}
Then for every morphism of schemes $g : S' \to S$ the pullback
$(g')^{-1}D$ is an effective Cartier divisor on $X' = S' \times_S X$
where $g' : X' \to X$ is the projection.
\end{lemma}

\begin{proof}
Using
Lemma \ref{lemma-characterize-effective-Cartier-divisor}
we translate this as follows into algebra. Let $A \to B$ be a ring
map and $h \in B$. Assume $h$ is a nonzerodivisor and that $B/hB$ is flat
over $A$. Then
$$
0 \to B \xrightarrow{h} B \to B/hB \to 0
$$
is a short exact sequence of $A$-modules with $B/hB$ flat over $A$. By
Algebra, Lemma \ref{algebra-lemma-flat-tor-zero}
this sequence remains exact on tensoring over $A$ with any module, in
particular with any $A$-algebra $A'$.
\end{proof}

\noindent
This lemma is the motivation for the following definition.

\begin{definition}
\label{definition-relative-effective-Cartier-divisor}
Let $f : X \to S$ be a morphism of schemes.
A {\it relative effective Cartier divisor} on $X/S$ is an
effective Cartier divisor $D \subset X$ such that $D \to S$
is a flat morphism of schemes.
\end{definition}

\noindent
We warn the reader that this may be nonstandard notation.
In particular, in \cite[IV, Section 21.15]{EGA} the notion of a
relative divisor is discussed only when $X \to S$ is flat and
locally of finite presentation. Our definition is a bit more general.
However, it turns out that if $x \in D$ then $X \to S$ is
flat at $x$ in many cases (but not always).

\begin{lemma}
\label{lemma-sum-relative-effective-Cartier-divisor}
Let $f : X \to S$ be a morphism of schemes. If $D_1, D_2 \subset X$
are relative effective Cartier divisor on $X/S$ then so
is $D_1 + D_2$ (Definition \ref{definition-sum-effective-Cartier-divisors}).
\end{lemma}

\begin{proof}
This translates into the following algebra fact:
Let $A \to B$ be a ring map and $h_1, h_2 \in B$.
Assume the $h_i$ are nonzerodivisors and that $B/h_iB$ is flat over $A$.
Then $h_1h_2$ is a nonzerodivisor and $B/h_1h_2B$ is flat over $A$.
The reason is that we have a short exact sequence
$$
0 \to B/h_1B \to B/h_1h_2B \to B/h_2B \to 0
$$
where the first arrow is given by multiplication by $h_2$. Since
the outer two are flat modules over $A$, so is the middle one, see
Algebra, Lemma \ref{algebra-lemma-flat-ses}.
\end{proof}

\begin{lemma}
\label{lemma-difference-relative-effective-Cartier-divisor}
Let $f : X \to S$ be a morphism of schemes. If $D_1, D_2 \subset X$
are relative effective Cartier divisor on $X/S$ and $D_1 \subset D_2$
as closed subschemes, then the effective Cartier divisor $D$
such that $D_2 = D_1 + D$
(Lemma \ref{lemma-difference-effective-Cartier-divisors}) is
a relative effective Cartier divisor on $X/S$.
\end{lemma}

\begin{proof}
This translates into the following algebra fact:
Let $A \to B$ be a ring map and $h_1, h_2 \in B$.
Assume the $h_i$ are nonzerodivisors, that $B/h_iB$ is flat over $A$, and
that $(h_2) \subset (h_1)$. Then we can write $h_2 = h h_1$
where $h \in B$ is a nonzerodivisor. We get a short exact sequence
$$
0 \to B/hB \to B/h_2B \to B/h_1B \to 0
$$
where the first arrow is given by multiplication by $h_1$. Since
the right two are flat modules over $A$, so is the middle one, see
Algebra, Lemma \ref{algebra-lemma-flat-ses}.
\end{proof}

\begin{lemma}
\label{lemma-flat-at-x}
Let $f : X \to S$ be a morphism of schemes.
Let $D \subset X$ be a relative effective Cartier divisor on $X/S$.
If $x \in D$ and $\mathcal{O}_{X, x}$ is Noetherian, then $f$ is flat at $x$.
\end{lemma}

\begin{proof}
Set $A = \mathcal{O}_{S, f(x)}$ and $B = \mathcal{O}_{X, x}$.
Let $h \in B$ be an element which generates the ideal of $D$.
Then $h$ is a nonzerodivisor in $B$ such that $B/hB$ is a flat
local $A$-algebra. Let $I \subset A$ be a finitely generated ideal.
Consider the commutative diagram
$$
\xymatrix{
0 \ar[r] &
B \ar[r]_h &
B \ar[r] &
B/hB \ar[r] & 0 \\
0 \ar[r] &
B \otimes_A I \ar[r]^h \ar[u] &
B \otimes_A I \ar[r] \ar[u] &
B/hB \otimes_A I \ar[r] \ar[u] & 0
}
$$
The lower sequence is short exact as $B/hB$ is flat over $A$, see
Algebra, Lemma \ref{algebra-lemma-flat-tor-zero}.
The right vertical arrow is injective as $B/hB$ is flat over $A$, see
Algebra, Lemma \ref{algebra-lemma-flat}.
Hence multiplication by $h$ is surjective on the kernel $K$ of
the middle vertical arrow. By Nakayama's lemma, see
Algebra, Lemma \ref{algebra-lemma-NAK}
we conclude that $K= 0$. Hence $B$ is flat over $A$, see
Algebra, Lemma \ref{algebra-lemma-flat}.
\end{proof}

\noindent
The following lemma relies on the algebraic version of
openness of the flat locus. The scheme theoretic version can be found in
More on Morphisms, Section \ref{more-morphisms-section-open-flat}.

\begin{lemma}
\label{lemma-flat-relative-Cartier-divisor}
Let $f : X \to S$ be a morphism of schemes.
Let $D \subset X$ be a relative effective Cartier divisor.
If $f$ is locally of finite presentation, then there exists
an open subscheme $U \subset X$ such that $D \subset U$ and
such that $f|_U : U \to S$ is flat.
\end{lemma}

\begin{proof}
Pick $x \in D$. It suffices to find an open neighbourhood $U \subset X$
of $x$ such that $f|_U$ is flat. Hence the lemma reduces to the case
that $X = \Spec(B)$ and $S = \Spec(A)$ are affine
and that $D$ is given by a nonzerodivisor $h \in B$. By assumption
$B$ is a finitely presented $A$-algebra and $B/hB$ is a flat
$A$-algebra. We are going to use absolute Noetherian approximation.

\medskip\noindent
Write $B = A[x_1, \ldots, x_n]/(g_1, \ldots, g_m)$. Assume
$h$ is the image of $h' \in A[x_1, \ldots, x_n]$. Choose a finite type
$\mathbf{Z}$-subalgebra $A_0 \subset A$ such that all the coefficients
of the polynomials $h', g_1, \ldots, g_m$ are in $A_0$. Then we can set
$B_0 = A_0[x_1, \ldots, x_n]/(g_1, \ldots, g_m)$ and $h_0$ the image
of $h'$ in $B_0$. Then $B = B_0 \otimes_{A_0} A$ and
$B/hB = B_0/h_0B_0 \otimes_{A_0} A$. By Algebra, Lemma
\ref{algebra-lemma-flat-finite-presentation-limit-flat}
we may, after enlarging $A_0$, assume that $B_0/h_0B_0$ is flat
over $A_0$. Let $K_0 = \Ker(h_0 : B_0 \to B_0)$.
As $B_0$ is of finite type over $\mathbf{Z}$ we see that $K_0$ is
a finitely generated ideal. Let $A_1 \subset A$ be a finite type
$\mathbf{Z}$-subalgebra containing $A_0$ and denote $B_1$, $h_1$, $K_1$
the corresponding objects over $A_1$. By
More on Algebra, Lemma \ref{more-algebra-lemma-base-change-H1-regular}
the map $K_0 \otimes_{A_0} A_1 \to K_1$ is surjective. On the other hand,
the kernel of $h : B \to B$ is zero by assumption. Hence every element
of $K_0$ maps to zero in $K_1$ for sufficiently large subrings
$A_1 \subset A$. Since $K_0$ is finitely generated, we conclude that
$K_1 = 0$ for a suitable choice of $A_1$.

\medskip\noindent
Set $f_1 : X_1 \to S_1$ equal to $\Spec$ of the
ring map $A_1 \to B_1$. Set $D_1 = \Spec(B_1/h_1B_1)$.
Since $B = B_1 \otimes_{A_1} A$, i.e., $X = X_1 \times_{S_1} S$,
it now suffices to prove the lemma for $X_1 \to S_1$ and the relative
effective Cartier divisor $D_1$, see
Morphisms, Lemma \ref{morphisms-lemma-base-change-module-flat}.
Hence we have reduced to the case where $A$ is a Noetherian ring.
In this case we know that the ring map $A \to B$ is flat at every
prime $\mathfrak q$ of $V(h)$ by
Lemma \ref{lemma-flat-at-x}.
Combined with the fact that the flat locus is open in this case, see
Algebra, Theorem \ref{algebra-theorem-openness-flatness}
we win.
\end{proof}

\noindent
There is also the following lemma (whose idea is apparently
due to Michael Artin, see \cite{Nobile}) which needs no finiteness
assumptions at all.

\begin{lemma}
\label{lemma-michael-artin}
Let $f : X \to S$ be a morphism of schemes.
Let $D \subset X$ be a relative effective Cartier divisor on $X/S$.
If $f$ is flat at all points of $X \setminus D$, then $f$ is flat.
\end{lemma}

\begin{proof}
This translates into the following algebra fact:
Let $A \to B$ be a ring map and $h \in B$.
Assume $h$ is a nonzerodivisor, that $B/hB$ is flat over $A$, and
that the localization $B_h$ is flat over $A$. Then $B$ is flat over $A$.
The reason is that we have a short exact sequence
$$
0 \to B \to B_h \to \colim_n (1/h^n)B/B \to 0
$$
and that the second and third terms are flat over $A$, which implies
that $B$ is flat over $A$ (see
Algebra, Lemma \ref{algebra-lemma-flat-ses}). Note that a filtered
colimit of flat modules is flat (see
Algebra, Lemma \ref{algebra-lemma-colimit-flat})
and that by induction on $n$ each $(1/h^n)B/B \cong B/h^nB$ is flat over
$A$ since it fits into the short exact sequence
$$
0 \to B/h^{n - 1}B \xrightarrow{h} B/h^nB \to B/hB \to 0
$$
Some details omitted.
\end{proof}

\begin{example}
\label{example-relative-cartier-ambient-space-not-flat}
Here is an example of a relative effective Cartier divisor $D$ where the
ambient scheme is not flat in a neighbourhood of $D$. Namely, let
$A = k[t]$ and
$$
B = k[t, x, y, x^{-1}y, x^{-2}y, \ldots]/(ty, tx^{-1}y, tx^{-2}y, \ldots)
$$
Then $B$ is not flat over $A$ but $B/xB \cong A$ is flat over $A$.
Moreover $x$ is a nonzerodivisor and hence defines a relative effective
Cartier divisor in $\Spec(B)$ over $\Spec(A)$.
\end{example}

\noindent
If the ambient scheme is flat and locally of finite presentation over
the base, then we can characterize a relative effective Cartier divisor
in terms of its fibres. See also
More on Morphisms, Lemma \ref{more-morphisms-lemma-slice-given-element}
for a slightly different take on this lemma.

\begin{lemma}
\label{lemma-fibre-Cartier}
Let $\varphi : X \to S$ be a flat morphism which is locally of finite
presentation. Let $Z \subset X$ be a closed subscheme.
Let $x \in Z$ with image $s \in S$.
\begin{enumerate}
\item If $Z_s \subset X_s$ is a Cartier divisor in a neighbourhood of $x$,
then there exists an open $U \subset X$ and a
relative effective Cartier divisor $D \subset U$ such that
$Z \cap U \subset D$ and $Z_s \cap U = D_s$.
\item If $Z_s \subset X_s$ is a Cartier divisor in a neighbourhood of $x$,
the morphism $Z \to X$ is of finite presentation, and $Z \to S$ is flat at
$x$, then we can choose $U$ and $D$ such that $Z \cap U = D$.
\item If $Z_s \subset X_s$ is a Cartier divisor in a neighbourhood of $x$
and $Z$ is a locally principal closed subscheme of $X$ in a neighbourhood
of $x$, then we can choose $U$ and $D$ such that $Z \cap U = D$.
\end{enumerate}
In particular, if $Z \to S$ is locally of finite presentation and flat and
all fibres $Z_s \subset X_s$ are effective Cartier divisors, then
$Z$ is a relative effective Cartier divisor. Similarly, if $Z$
is a locally principal closed subscheme of $X$ such that all fibres
$Z_s \subset X_s$ are effective Cartier divisors, then
$Z$ is a relative effective Cartier divisor.
\end{lemma}

\begin{proof}
Choose affine open neighbourhoods $\Spec(A)$ of $s$ and
$\Spec(B)$ of $x$ such that
$\varphi(\Spec(B)) \subset \Spec(A)$.
Let $\mathfrak p \subset A$ be the prime ideal corresponding to $s$.
Let $\mathfrak q \subset B$ be the prime ideal corresponding to $x$.
Let $I \subset B$ be the ideal corresponding to $Z$.
By the initial assumption of the lemma we know that
$A \to B$ is flat and of finite presentation.
The assumption in (1) means that, after shrinking $\Spec(B)$, we may
assume $I(B \otimes_A \kappa(\mathfrak p))$ is generated by a single
element which is a nonzerodivisor in $B \otimes_A \kappa(\mathfrak p)$.
Say $f \in I$ maps to this generator. We claim that after inverting
an element $g \in B$, $g \not \in \mathfrak q$ the closed subscheme
$D = V(f) \subset \Spec(B_g)$ is a relative effective Cartier
divisor.

\medskip\noindent
By
Algebra, Lemma \ref{algebra-lemma-flat-finite-presentation-limit-flat}
we can find a flat finite type ring map $A_0 \to B_0$ of Noetherian
rings, an element $f_0 \in B_0$, a ring map $A_0 \to A$ and an isomorphism
$A \otimes_{A_0} B_0 \cong B$. If $\mathfrak p_0 = A_0 \cap \mathfrak p$
then we see that
$$
B \otimes_A \kappa(\mathfrak p) =
\left(B_0 \otimes_{A_0} \kappa(\mathfrak p_0)\right)
\otimes_{\kappa(\mathfrak p_0))} \kappa(\mathfrak p)
$$
hence $f_0$ is a nonzerodivisor in $B_0 \otimes_{A_0} \kappa(\mathfrak p_0)$.
By
Algebra, Lemma \ref{algebra-lemma-grothendieck}
we see that $f_0$ is a nonzerodivisor in $(B_0)_{\mathfrak q_0}$
where $\mathfrak q_0 = B_0 \cap \mathfrak q$ and
that $(B_0/f_0B_0)_{\mathfrak q_0}$ is flat over $A_0$. Hence by
Algebra, Lemma \ref{algebra-lemma-regular-sequence-in-neighbourhood}
and
Algebra, Theorem \ref{algebra-theorem-openness-flatness}
there exists a $g_0 \in B_0$, $g_0 \not \in \mathfrak q_0$ such
that $f_0$ is a nonzerodivisor in $(B_0)_{g_0}$ and such that
$(B_0/f_0B_0)_{g_0}$ is flat over $A_0$. Hence we see that
$D_0 = V(f_0) \subset \Spec((B_0)_{g_0})$ is a relative effective
Cartier divisor. Since we know that this property is preserved under
base change, see
Lemma \ref{lemma-relative-Cartier},
we obtain the claim mentioned above with $g$ equal to the image of $g_0$
in $B$.

\medskip\noindent
At this point we have proved (1). To see (2) consider the closed
immersion $Z \to D$. The surjective ring map
$u : \mathcal{O}_{D, x} \to \mathcal{O}_{Z, x}$
is a map of flat local $\mathcal{O}_{S, s}$-algebras which
are essentially of finite presentation, and which becomes an
isomorphisms after dividing by $\mathfrak m_s$. Hence it is
an isomorphism, see
Algebra, Lemma \ref{algebra-lemma-mod-injective-general}.
It follows that $Z \to D$ is an isomorphism in a neighbourhood
of $x$, see
Algebra, Lemma \ref{algebra-lemma-local-isomorphism}.
To see (3), after possibly shrinking $U$ we may assume that
the ideal of $D$ is generated by a single nonzerodivisor $f$
and the ideal of $Z$ is generated by an element $g$. Then
$f = gh$. But $g|_{U_s}$ and $f|_{U_s}$ cut out the same
effective Cartier divisor in a neighbourhood of $x$. Hence
$h|_{X_s}$ is a unit in $\mathcal{O}_{X_s, x}$, hence $h$ is
a unit in $\mathcal{O}_{X, x}$ hence $h$ is a unit in an
open neighbourhood of $x$. I.e., $Z \cap U = D$ after shrinking $U$.

\medskip\noindent
The final statements of the lemma follow immediately from
parts (2) and (3), combined with the fact that $Z \to S$
is locally of finite presentation if and only if $Z \to X$ is
of finite presentation, see
Morphisms, Lemmas \ref{morphisms-lemma-composition-finite-presentation} and
\ref{morphisms-lemma-finite-presentation-permanence}.
\end{proof}



